\chapter{Conjunto de datos} \label{chap:dataset}

En este capítulo se describe el origen, estructura y preprocesamiento del conjunto de datos utilizado en esta tesis. El corpus proviene de la tarea compartida REST-MEX 2025, organizada en el marco de IberLEF 2025. Se detalla el contexto del evento, la composición del dataset original, y el proceso de filtrado aplicado para obtener una matriz de interacciones con densidad suficiente para el entrenamiento de sistemas de recomendación.

%% ============================================================
%% EVENTO REST-MEX
%% ============================================================
\section{La tarea compartida REST-MEX} \label{sec:restmex}

REST-MEX (\textit{Researching Sentiment Evaluation in Text for Mexican Magical Towns}) es una tarea compartida organizada dentro del foro de evaluación IberLEF (\textit{Iberian Languages Evaluation Forum}), cuyo objetivo es establecer referencias comparativas (\textit{benchmarks}) para sistemas de procesamiento de lenguaje natural en el ámbito del turismo mexicano \citep{alvarez-carmona_overview_2025}.

La tarea se lanzó en 2021 con un enfoque en la predicción de polaridad de reseñas turísticas en español \citep{alvarez-carmona_overview_2021}. En ediciones subsecuentes, se incorporó la clasificación del tipo de servicio reseñado (hotel, restaurante o atractivo) \citep{alvarez-carmona_overview_2022} y la identificación del país de origen de la reseña \citep{alvarez-carmona_overview_2023}. La edición 2025 marca la cuarta iteración del evento e introduce cambios significativos: (i)~la clasificación geográfica se refina a nivel de Pueblo Mágico específico (40 pueblos predefinidos), (ii)~el corpus se amplía considerablemente, y (iii)~se fomenta el uso de técnicas avanzadas como ingeniería de \textit{prompts}, binarización de clases y selección de instancias basada en centroides.

En esta edición, 32 equipos participantes enviaron un total de 70 sistemas válidos, explorando una amplia variedad de enfoques que incluyen modelos basados en Transformers, aprendizaje por transferencia, esquemas de binarización y estrategias de selección de instancias \citep{alvarez-carmona_overview_2025}. La tarea se estructura en tres subtareas simultáneas que, dada una reseña turística, requieren determinar:

\begin{enumerate}
    \item La \textbf{polaridad} del sentimiento, en una escala ordinal de 1 (muy malo) a 5 (muy bueno).
    \item El \textbf{tipo de servicio} reseñado: \textit{Attractive}, \textit{Hotel} o \textit{Restaurant}.
    \item El \textbf{Pueblo Mágico} específico al que refiere la reseña, seleccionado de una lista predefinida de 40 destinos.
\end{enumerate}

Si bien la Tarea 4 del evento REST-MEX contempla la construcción de un sistema de recomendación para Pueblos Mágicos, en esta tesis se extiende dicho planteamiento incorporando análisis de sentimientos basado en aspectos y optimización metaheurística para la fusión de componentes, aspectos que no están contemplados en la formulación original del evento.

%% ============================================================
%% CORPUS ORIGINAL
%% ============================================================
\section{Corpus REST-MEX 2025} \label{sec:corpus}

El corpus REST-MEX 2025 consta de 297,217 escritas por turistas que visitaron destinos representativos de México y compartieron sus opiniones en la plataforma TripAdvisor entre 2002 y 2024 \citep{alvarez-carmona_overview_2025}. Cada reseña está etiquetada con una calificación de polaridad en una escala ordinal de cinco puntos $\{1, 2, 3, 4, 5\}$, el tipo de negocio (Attractive, Hotel, Restaurant) y el Pueblo Mágico correspondiente.

El corpus presenta un notable desbalance de clases en los tres ejes de clasificación. En el eje de polaridad, la clase 5 (muy bueno) concentra el $65.6\%$ de las instancias de entrenamiento ($136{,}561$ de $208{,}051$), mientras que las clases 1 y 2 representan conjuntamente menos del $5.3\%$. En el eje geográfico, tres pueblos (Tulum, Isla Mujeres y San Cristóbal de las Casas) concentran el $42.2\%$ de las instancias.

Para los fines de esta tesis, se recopiló un conjunto extendido de reseñas a partir de los archivos fuente del evento. El proceso de carga concatenó 1,559 archivos CSV, resultando en un dataset crudo de 313,836 registros con 12 columnas originales.

%% ============================================================
%% ESQUEMA Y LIMPIEZA
%% ============================================================
\section{Estructura y limpieza del dataset} \label{sec:cleaning}

Cada registro del dataset crudo contiene los siguientes campos:

\begin{table}[ht]
\centering
\caption{Esquema del dataset crudo de reseñas turísticas.}
\label{tab:schema}
\begin{tabular}{lll}
\toprule
\textbf{Campo} & \textbf{Tipo} & \textbf{Descripción} \\
\midrule
Author        & Texto     & Seudónimo del usuario en TripAdvisor \\
Titulo        & Texto     & Título de la reseña \\
Review        & Texto     & Cuerpo de la reseña \\
Calificacion  & Númerico  & Calificación numérica (1--5) \\
OrigenAutor   & Texto     & País de origen del autor \\
FechaOpinion  & Fecha     & Fecha de publicación de la reseña \\
FechaEstadia  & Fecha     & Mes y año de la estancia \\
Pueblo        & Texto     & Pueblo Mágico del establecimiento \\
Estado        & Texto     & Estado de la república \\
Pais          & Texto     & País del establecimiento \\
Tipo          & Texto     & Tipo de servicio (Hotel/Restaurant/Attractive) \\
Lugar         & Texto     & Nombre del establecimiento \\
\bottomrule
\end{tabular}
\end{table}

El proceso de limpieza consistió en:

\begin{enumerate}
    \item \textbf{Eliminación de columnas redundantes}: se descartaron \texttt{OrigenAutor} (53.5\% de valores nulos), \texttt{FechaOpinion} (redundante con \texttt{FechaEstadia}) y \texttt{Pais} (valor constante ``Mexico'').
    \item \textbf{Eliminación de registros incompletos}: se eliminaron filas con valores nulos en los campos \texttt{Author}, \texttt{Calificacion}, \texttt{Review} o \texttt{FechaEstadia}, reduciendo el dataset de $313{,}836$ a \textbf{306,854 registros}.
    \item \textbf{Conversión de fechas}: el campo \texttt{FechaEstadia}, originalmente en formato textual en español (e.g., ``enero de 2023''), se convirtió a formato \texttt{datetime} mediante un mapeo de nombres de meses en español a sus equivalentes numéricos.
\end{enumerate}

El dataset limpio resultante contiene 9 columnas y abarca estancias desde 2001 hasta 2023.

%% ============================================================
%% FILTRADO DE DENSIDAD
%% ============================================================
\section{Filtrado para densidad de interacciones} \label{sec:filtering}

Para el entrenamiento de algoritmos de filtrado colaborativo, es necesario que la matriz de interacciones usuario-ítem presente una densidad mínima que permita calcular similitudes confiables entre usuarios o entre ítems. El dataset limpio de $306{,}854$ reseñas contiene un gran número de usuarios que han reseñado establecimientos en un solo Pueblo Mágico, lo cual resulta insuficiente para la tarea de recomendación a nivel de pueblos.

Se diseñó un procedimiento de filtrado iterativo con dos criterios de densidad:

\begin{enumerate}
    \item \textbf{Usuarios activos}: retener únicamente usuarios que hayan reseñado establecimientos en al menos $\tau_u = 3$ Pueblos Mágicos distintos.
    \item \textbf{Pueblos con cobertura}: retener únicamente Pueblos Mágicos que contengan al menos $\tau_p = 10$ establecimientos (\textit{Lugares}) distintos con reseñas.
\end{enumerate}

Dado que la eliminación de usuarios puede dejar pueblos sin cobertura suficiente (y viceversa), el filtrado se aplica de forma iterativa hasta convergencia, es decir, hasta que ningún registro adicional sea eliminado. El Algoritmo~\ref{alg:filter} formaliza este procedimiento.

\begin{algorithm}[ht]
\caption{Filtrado iterativo para densidad de interacciones.}
\label{alg:filter}
\begin{algorithmic}[1]
\REQUIRE Dataset $\mathcal{D}$, umbrales $\tau_u$, $\tau_p$
\ENSURE Dataset filtrado $\mathcal{D}'$
\STATE $\mathcal{D}' \leftarrow \mathcal{D}$
\REPEAT
    \STATE $n_{\text{prev}} \leftarrow |\mathcal{D}'|$
    \STATE $\mathcal{U}_{\text{dense}} \leftarrow \{u \in \mathcal{U} \mid |\{p : \exists\, r \in \mathcal{D}' \text{ con } r.\text{Author} = u,\; r.\text{Pueblo} = p\}| \geq \tau_u\}$
    \STATE $\mathcal{D}' \leftarrow \{r \in \mathcal{D}' \mid r.\text{Author} \in \mathcal{U}_{\text{dense}}\}$
    \STATE $\mathcal{P}_{\text{dense}} \leftarrow \{p \in \mathcal{P} \mid |\{l : \exists\, r \in \mathcal{D}' \text{ con } r.\text{Pueblo} = p,\; r.\text{Lugar} = l\}| \geq \tau_p\}$
    \STATE $\mathcal{D}' \leftarrow \{r \in \mathcal{D}' \mid r.\text{Pueblo} \in \mathcal{P}_{\text{dense}}\}$
\UNTIL{$|\mathcal{D}'| = n_{\text{prev}}$}
\RETURN $\mathcal{D}'$
\end{algorithmic}
\end{algorithm}

%% ============================================================
%% DATASET RESULTANTE
%% ============================================================
\section{Dataset de trabajo} \label{sec:working_dataset}

La aplicación del filtrado iterativo con $\tau_u = 3$ y $\tau_p = 10$ reduce el dataset de $306{,}854$ a \textbf{79,264 registros}, correspondientes a \textbf{9,582 usuarios} y \textbf{48 Pueblos Mágicos}. La Tabla~\ref{tab:dataset_stats} resume las estadísticas principales antes y después del filtrado.

\begin{table}[ht]
\centering
\caption{Estadísticas del dataset antes y después del filtrado de densidad.}
\label{tab:dataset_stats}
\begin{tabular}{lrr}
\toprule
\textbf{Métrica} & \textbf{Dataset limpio} & \textbf{Dataset filtrado} \\
\midrule
Reseñas totales    & 306,854 & 79,264 \\
Usuarios únicos    & ---     & 9,582  \\
Pueblos Mágicos    & ---     & 48     \\
Mín. pueblos/usuario & 1     & 3      \\
Máx. pueblos/usuario & ---   & 28     \\
Mín. lugares/pueblo  & ---   & 10     \\
\bottomrule
\end{tabular}
\end{table}

La distribución de pueblos en el dataset filtrado presenta un marcado desbalance geográfico. Tulum es el pueblo con mayor representación ($10{,}051$ títulos únicos), seguido de Isla Mujeres ($6{,}082$) y San Cristóbal de las Casas ($4{,}902$), mientras que pueblos como Tecate ($175$), San Sebastián del Oeste ($193$) y Santiago ($242$) tienen una representación considerablemente menor.

En términos de la matriz de utilidad para el sistema de recomendación---donde las filas representan usuarios y las columnas Pueblos Mágicos---la dispersión (\textit{sparsity}) se calcula como:
\begin{equation} \label{eq:sparsity}
    \text{sparsity} = 1 - \frac{|\{(u, p) : \exists\, r \in \mathcal{D}'\}|}{|\mathcal{U}| \times |\mathcal{P}|} = 1 - \frac{n_{\text{interacciones}}}{9{,}582 \times 48} \approx 90.32\%
\end{equation}

Esta dispersión del $90.32\%$ implica que cada usuario ha interactuado, en promedio, con menos de 5 de los 48 pueblos disponibles. En cuanto a las calificaciones, la media global es $\bar{r} = 4.41$ con el $61.3\%$ de las calificaciones siendo de 5 estrellas, reflejando el sesgo positivo característico de las plataformas de reseñas turísticas. Estos desafíos---alta dispersión y desbalance en calificaciones---motivan la adopción de un enfoque híbrido que combine señales colaborativas con información semántica extraída de las reseñas textuales, como se detalla en los capítulos subsecuentes.
